\documentclass[11pt, a4paper]{article}

\usepackage[margin=2.5cm]{geometry}
\usepackage{booktabs}
\usepackage{longtable}
\usepackage{array}
\usepackage{tabularx}
\usepackage[table]{xcolor}
\usepackage{hyperref}
\usepackage{titlesec}
\usepackage{enumitem}
\usepackage{fancyhdr}
\usepackage{amsmath}
\usepackage{amssymb}
\usepackage{microtype}
\usepackage{parskip}
\usepackage{mdframed}
\usepackage{tikz}
\usetikzlibrary{shapes.geometric, arrows.meta, positioning, fit, backgrounds}

% --- colour palette (phosphor green theme) ---
\definecolor{phosphor}{RGB}{20, 200, 20}
\definecolor{dimgreen}{RGB}{10, 120, 10}
\definecolor{darkbg}{RGB}{8, 20, 8}
\definecolor{accentcyan}{RGB}{0, 180, 220}
\definecolor{accentred}{RGB}{220, 40, 20}
\definecolor{tabhead}{RGB}{20, 60, 20}
\definecolor{rowalt}{RGB}{230, 245, 230}
\definecolor{notebg}{RGB}{240, 250, 240}

% --- hyperref ---
\hypersetup{
  colorlinks = true,
  linkcolor  = dimgreen,
  urlcolor   = accentcyan,
  citecolor  = dimgreen,
}

% --- header / footer ---
\pagestyle{fancy}
\fancyhf{}
\fancyhead[L]{\small\textcolor{dimgreen}{OAMP — Development Roadmap}}
\fancyhead[R]{\small\textcolor{dimgreen}{Phases 2--5}}
\fancyfoot[C]{\small\textcolor{dimgreen}{\thepage}}
\renewcommand{\headrulewidth}{0.4pt}
\renewcommand{\headrule}{\hbox to\headwidth{\color{dimgreen}\leaders\hrule height \headrulewidth\hfill}}

% --- section styling ---
\titleformat{\section}{\large\bfseries\color{dimgreen}}{}{0em}{}[\textcolor{dimgreen}{\titlerule[0.6pt]}]
\titleformat{\subsection}{\normalsize\bfseries\color{dimgreen!80!black}}{}{0em}{}
\titleformat{\subsubsection}{\normalsize\itshape\color{dimgreen!70!black}}{}{0em}{}

% --- note box ---
\newmdenv[
  backgroundcolor=notebg,
  linecolor=dimgreen,
  linewidth=1pt,
  leftmargin=0pt, rightmargin=0pt,
  innerleftmargin=10pt, innerrightmargin=10pt,
  innertopmargin=6pt, innerbottommargin=6pt,
]{notebox}

% --- table helpers ---
\newcolumntype{L}[1]{>{\raggedright\arraybackslash}p{#1}}
\newcolumntype{C}[1]{>{\centering\arraybackslash}p{#1}}

\begin{document}

% ============================================================
%  TITLE
% ============================================================
\begin{titlepage}
  \pagecolor{darkbg}
  \color{phosphor}
  \vspace*{3cm}
  \begin{center}
    {\Huge\bfseries Open Source Astrodynamic\\[0.3em]
     Mission Planning}\\[1.5em]
    {\LARGE\color{accentcyan} Development Roadmap}\\[0.8em]
    {\large Phases 2 \& 3}\\[3cm]
    \textcolor{phosphor!60!black}{\rule{0.5\textwidth}{0.4pt}}\\[1em]
    {\normalsize
      Precision Newtonian Mission Design \quad\textbullet\quad
      Post-Newtonian \& General Relativity}\\[2em]
    {\small\color{phosphor!60!black}
      Apache-2.0 \quad|\quad HortDog \quad|\quad \today}
  \end{center}
\end{titlepage}

\pagecolor{white}
\color{black}

% ============================================================
%  TABLE OF CONTENTS
% ============================================================
\tableofcontents
\newpage

% ============================================================
%  EXECUTIVE SUMMARY
% ============================================================
\section{Executive Summary}

Phase~1 delivered the project skeleton: a Newtonian $+$~J2 propagator with
DOP853 integration, an open-loop launch simulator, a SPICE kernel wrapper with
reproducibility-locked manifests, and a WebGPU~3D viewer showing three orbital
scenarios. All 10 unit tests pass; the stack is clean (ruff, TypeScript strict).

\textbf{Phase~2} keeps the physics Newtonian but raises the fidelity to the
level needed for real mission design: higher-order gravity, atmospheric drag,
solar radiation pressure, third-body perturbations, TLE ingest, $\Delta v$
manoeuvre planning, a Lambert/CasADi trajectory optimizer, live WebSocket
streaming, and a full manoeuvre-editor frontend.

\textbf{Phase~3} (\emph{shipped}) brings every backend capability to the
sandbox frontend (J2--J6, drag with MSIS, SRP, third-body Moon/Sun, finite
burns, four solver panels covering Hohmann/Lambert/multi-burn NLP/TLE), adds
a central-body selector with full Moon support (canonical $\mu$, $R$, $J_n$,
$\omega$), wires SPICE ephemerides into the renderer as a secondary-body
wireframe, and ties it together with an end-to-end cislunar TLI demo
(Lambert + J2 + third-body Moon + finite-burn-ready spacecraft).

\textbf{Phase~4} extends the dynamics layer with the circular-restricted
three-body problem (CR3BP), Lagrange-point computation, periodic-orbit
families with differential correction, invariant manifold tubes, and
rotating-frame visualization --- the foundations for low-energy / weak
stability boundary (WSB) transfers and ballistic capture into lunar orbit.

\textbf{Phase~5} adds post-Newtonian (PN) corrections and a full General
Relativity (GR) geodesic integrator on the DSX/Brumberg--Kopejkin BCRS metric,
with tetrad thrust modelling and a Ray-backed spacetime metric pre-bake cache
for performance.

\begin{notebox}
\textbf{Convention lock (unchanged from Phase~1):}
UTC ISO-8601 on the wire; TDB seconds since J2000 (= SPICE ET) internally.
All SI units in the engine. Apache-2.0 licence.
\end{notebox}

% ============================================================
%  PHASE 2
% ============================================================
\newpage
\section{Phase 2 — Precision Newtonian \& Mission Design}

\subsection{Goals}

A mission designer should be able to: load a TLE or specify an initial state,
apply atmospheric drag and solar radiation pressure, add impulsive $\Delta v$
burns, solve a Lambert transfer or run a fuel-optimal CasADi optimisation, and
watch the result stream live in the 3D viewer.

\subsection{2a — Dynamics Fidelity}

\subsubsection{Higher-Order Gravity}

Extend \texttt{oamp.dynamics.newtonian} with zonal harmonics $J_3$--$J_6$
(Vallado §8.7). The acceleration from $J_n$ is:

\begin{equation}
  \mathbf{a}_{J_n} = -\frac{\mu}{r^2} \frac{n+1}{2}
    \left(\frac{R_\oplus}{r}\right)^n J_n \,
    \nabla P_n(\sin\phi)
\end{equation}

where $P_n$ is the Legendre polynomial and $\phi$ the geocentric latitude.
Tesseral ($C_{nm}$, $S_{nm}$) terms are planned but deferred to Phase~2b+.

\subsubsection{Atmospheric Drag}

Replace the exponential density model in \texttt{launch.py} with
NRLMSISE-00 via the \texttt{nrlmsise00} Python package. The drag acceleration
is:

\begin{equation}
  \mathbf{a}_\text{drag} = -\tfrac{1}{2}\, \rho(h)\, C_D\, \frac{A}{m}\,
    |\mathbf{v}_\text{rel}|\,\mathbf{v}_\text{rel}
\end{equation}

where $\mathbf{v}_\text{rel}$ accounts for Earth's rotation
($\boldsymbol{\omega}_\oplus \times \mathbf{r}$). Vehicle cross-section $A$
and $C_D$ are added to the propagation request body.

\subsubsection{Solar Radiation Pressure}

Cannonball model with cylindrical shadow cones for Earth and Moon:

\begin{equation}
  \mathbf{a}_\text{SRP} = \nu \frac{P_\odot}{c}
    \frac{A_\text{eff}}{m} \hat{\mathbf{r}}_{\odot}
\end{equation}

where $\nu \in [0,1]$ is the shadow function (computed via SPICE body
positions already available) and $P_\odot = 4.56\times10^{-6}$~N~m$^{-2}$
at 1~AU.

\subsubsection{Third-Body Gravity}

Moon and Sun perturbations using SPICE \texttt{body\_state()} (already
implemented) with point-mass attractions and the indirect term:

\begin{equation}
  \mathbf{a}_\text{3b} = \sum_k \mu_k
    \left(\frac{\mathbf{r}_k - \mathbf{r}}{|\mathbf{r}_k - \mathbf{r}|^3}
          - \frac{\mathbf{r}_k}{|\mathbf{r}_k|^3}\right)
\end{equation}

\subsubsection{Symplectic Integrator}

Add SABA2 (symplectic Birkhoff--Adams) alongside DOP853 for long-arc
propagation where energy conservation over months matters more than step-size
adaptivity. Selectable via \texttt{integrator: "dop853" | "saba2"} in the
request.

\subsubsection{New API Flags}

\begin{center}
\begin{tabular}{L{3cm} L{4cm} L{5cm}}
  \toprule
  \textbf{Field} & \textbf{Type} & \textbf{Effect} \\
  \midrule
  \texttt{drag}       & \texttt{bool}        & NRLMSISE-00 drag \\
  \texttt{srp}        & \texttt{bool}        & Solar radiation pressure \\
  \texttt{third\_body}& \texttt{list[str]}   & Bodies e.g.\ \texttt{["moon","sun"]} \\
  \texttt{jn\_max}    & \texttt{int}         & Max zonal harmonic (2–6) \\
  \texttt{integrator} & \texttt{str}         & \texttt{"dop853"} or \texttt{"saba2"} \\
  \texttt{vehicle}    & \texttt{VehicleModel}& $m$, $A$, $C_D$, $C_R$ \\
  \bottomrule
\end{tabular}
\end{center}

\subsubsection{Tests}

\begin{itemize}[noitemsep]
  \item Energy drift $< 10^{-9}$ relative over 30-day arc (SABA2)
  \item Nodal regression matches NRLMSISE drag prediction within 5\%
  \item SRP produces correct signature in inclination evolution
  \item Moon third-body matches SPICE-derived perturbation to $< 1$~m/day
\end{itemize}

\subsection{2b — TLE Ingest}

Add \texttt{GET /tle?norad=\{id\}} which:
\begin{enumerate}[noitemsep]
  \item Fetches the current TLE from Celestrak
  \item Parses with \texttt{sgp4} (already evaluable from \texttt{conda-forge})
  \item Propagates to epoch via SGP4
  \item Converts mean-element state to osculating Cartesian (Brouwer transform)
  \item Returns an \texttt{InitialState} compatible with \texttt{POST /propagate}
\end{enumerate}

A manifest of pinned test TLEs (NORAD ID + epoch + expected state checksum)
mirrors the SPICE kernel manifest pattern for reproducible CI.

\subsection{2c — $\Delta v$ Manoeuvre Engine}

\subsubsection{Data Model}

\begin{verbatim}
class Maneuver(BaseModel):
    t_tdb: float                # TDB seconds since J2000
    dv_ric: list[float]         # [radial, in-track, cross-track], m/s
\end{verbatim}

\texttt{POST /propagate} accepts an optional \texttt{maneuvers: list[Maneuver]};
the integrator applies instantaneous velocity kicks at each epoch by
transforming RIC $\to$ ECI and restarting the ODE.

\subsubsection{Analytical Solvers}

\begin{center}
\begin{tabular}{L{3.5cm} L{9cm}}
  \toprule
  \textbf{Endpoint} & \textbf{Description} \\
  \midrule
  \texttt{POST /optimize/hohmann}  &
    Closed-form two-impulse Hohmann transfer. Returns $\Delta v_1$,
    $\Delta v_2$, transfer time. \\[4pt]
  \texttt{POST /optimize/lambert}  &
    Universal-variable Lambert solver (Izzo/Gooding algorithm) for
    arbitrary two-point boundary-value transfer. \\
  \bottomrule
\end{tabular}
\end{center}

\subsubsection{Numerical Optimiser}

\texttt{POST /optimize/trajectory} uses CasADi (already in the \texttt{default}
pixi environment) to solve the fuel-optimal NLP:

\begin{equation}
  \min_{\{\Delta \mathbf{v}_k\}} \sum_k |\Delta \mathbf{v}_k|
  \quad \text{s.t.} \quad
  \mathbf{x}_{k+1} = \phi(\mathbf{x}_k + \Delta\mathbf{v}_k,\, \Delta t_k),\;
  \mathbf{x}_N = \mathbf{x}_\text{target}
\end{equation}

The propagation $\phi$ is the DOP853 integrator wrapped as a CasADi callback.
Warm-starting from the Lambert solution typically converges in $< 30$~iterations.

\subsection{2d — Real-time WebSocket Streaming}

Replace the echo-only \texttt{/ws} channel with a streaming propagator. The
protocol uses msgpack (already a dependency):

\begin{enumerate}[noitemsep]
  \item Client sends \texttt{\{scenario\_id, t\_start, t\_end, dt, config\}}
  \item Server propagates in chunks of $N$ steps on a Ray actor (or asyncio for
        small jobs)
  \item Each chunk is msgpack-encoded and sent as a binary WebSocket frame
  \item A terminal \texttt{\{done: true, elapsed\_ms\}} frame closes the stream
\end{enumerate}

This enables the frontend time-scrubber to show live progress for long
multi-day propagations.

\subsection{2e — Frontend}

\begin{center}
\begin{tabularx}{\textwidth}{L{3.5cm} X}
  \toprule
  \textbf{Feature} & \textbf{Detail} \\
  \midrule
  Manoeuvre editor &
    Click on trajectory to insert a burn node; drag the arrow glyph to adjust
    $\Delta v$ magnitude and direction in the RIC frame. \\[4pt]
  Time scrubber &
    Slider replays the trajectory; a spacecraft marker moves along the path.
    Shows $t_\text{UTC}$, $r$, $v$, and remaining $\Delta v$ budget. \\[4pt]
  Multi-trajectory overlay &
    Compare baseline vs optimised trajectories side-by-side with per-trajectory
    colour coding and a difference-norm plot in a side panel. \\[4pt]
  Velocity / thrust vectors &
    Short arrow glyphs at the spacecraft position; length proportional to
    $|\mathbf{v}|$ or $|\mathbf{f}|$ with configurable scale. \\[4pt]
  WebGL2 fallback &
    2D ground-track canvas (equirectangular) for browsers without WebGPU.
    Uses the same trajectory data from the REST API. \\[4pt]
  Scenario save/load &
    JSON export of full mission state (initial conditions, manoeuvres,
    vehicle model, propagation config). Drag-and-drop import. \\
  \bottomrule
\end{tabularx}
\end{center}

\subsection{Phase 2 Sequencing}

\begin{center}
\begin{tabular}{L{5cm} C{2.5cm} L{5.5cm}}
  \toprule
  \textbf{Work item} & \textbf{Est.\ duration} & \textbf{Dependency} \\
  \midrule
  2a Drag + third-body     & 2--3 weeks & Phase 1 \\
  2b TLE ingest            & 1 week     & Phase 1 \\
  2c Lambert + CasADi NLP  & 2--3 weeks & 2a \\
  2d WebSocket streaming   & 1 week     & Phase 1 \\
  2e Frontend              & 3--4 weeks & 2c, 2d (parallel) \\
  \midrule
  \textbf{Total}           & \textbf{9--12 weeks} & \\
  \bottomrule
\end{tabular}
\end{center}

% ============================================================
%  PHASE 3 — Sandbox Integration & Lunar System (SHIPPED)
% ============================================================
\newpage
\section{Phase 3 — Sandbox Integration \& Lunar System}

\subsection{Goals}

Make every backend capability visible and editable in the WebGPU sandbox,
extend the central-body abstraction to the Moon, and use SPICE ephemerides
to render the Moon at its inertial position --- producing an end-to-end
cislunar transfer demo that exercises the full Phase~2 stack
(higher-order gravity, third-body perturbation, finite burns, Lambert).

\subsection{3.1 — Perturbation editor expansion (shipped)}

The maneuver editor sidebar gained a full Perturbations section exposing
every flag the backend already accepts:

\begin{itemize}[leftmargin=*]
  \item Zonal harmonics dropdown ($J_2$\,\dots\,$J_6$) bound to
        \texttt{PropagateRequest.jn\_max}
  \item Atmospheric drag with model selector
        (\texttt{exponential} \,|\, \texttt{msis}); vehicle inputs share
        a single section
  \item Solar radiation pressure (cannonball) with $C_r$, area
  \item Third-body checkboxes for Moon and Sun
        (forwarded as \texttt{third\_body: ["MOON", "SUN"]})
  \item Live ``active forces'' readout displaying what the server actually
        applied (\texttt{response.perturbations[]})
\end{itemize}

\subsection{3.2 — Solver panel (shipped)}

A single sidebar section with a solver-type dropdown and a dynamically
rendered form:

\begin{center}
\rowcolors{2}{rowalt}{white}
\begin{tabularx}{\textwidth}{>{\bfseries}l X X}
  \rowcolor{tabhead}\color{white}\textbf{Solver}
    & \color{white}\textbf{Inputs}
    & \color{white}\textbf{Loads into editor} \\
  Hohmann       & $r_1$, $r_2$, $\mu$
                & 2 impulsive maneuvers, IC at LEO \\
  Lambert       & $\vec r_1$, $\vec r_2$, TOF
                & IC = $(\vec r_1, \vec v_1^{\text{Lambert}})$, duration $=$ TOF \\
  Multi-burn NLP & $N$ epochs, target $\vec r_f, \vec v_f$
                & Δv list + convergence stats (inertial; displayed) \\
  TLE           & NORAD ID \emph{or} line1/line2 (+ \texttt{at\_utc})
                & ECI state at epoch, period auto-set \\
\end{tabularx}
\end{center}

Each solver writes results back to \texttt{editorState} and auto-propagates,
so the user sees the orbit immediately.

\subsection{3.3 — Central-body selector + lunar constants}

\texttt{PropagateRequest.body\_name: "EARTH" | "MOON" | "SUN"} resolves to
the canonical \texttt{Body} constant (correct $\mu$, $R$, $J_n$, $\omega$);
the legacy \texttt{mu}/\texttt{body\_radius} fields are kept for back-compat.
Lunar zonal harmonics $J_2$\,\dots\,$J_6$ are populated from the GRGM1200A
low-order solution. The frontend dropdown repopulates the initial state
with a sensible low circular orbit at the body's default altitude and
clears drag/SRP (which require an atmosphere).

\subsection{3.4 — Moon rendering}

The WebGPU renderer gained \texttt{setSecondaryBody(position, radius, color)}:
a low-resolution wireframe sphere translated to the inertial Moon position.
A new endpoint \texttt{POST /spice/ephemeris} accepts a list of TDB-second
offsets from a UTC epoch and returns the corresponding body states ---
batching what \texttt{/spice/state} did for a single time. The frontend
samples Moon positions at $\sim$60 points across the active trajectory and
advances the Moon as the time scrubber moves (binary-search lookup).

\subsection{3.5 — Cislunar TLI demo}

A top-bar button that ties the whole stack together:

\begin{enumerate}[leftmargin=*]
  \item Query Moon position at the demo epoch from SPICE
  \item Solve Lambert: $200\,$km LEO $\to$ Moon position in $3.5\,$d
  \item Load TLI Δv as an impulsive RIC maneuver
  \item Propagate with $J_2$ + third-body Moon
  \item Render trajectory + Moon wireframe + apse markers
\end{enumerate}

Graceful degradation: if SPICE kernels aren't loaded the status bar prompts
the user to run \texttt{pixi run kernels} rather than failing silently.

\subsection{3.6 — Polish}

\begin{itemize}[leftmargin=*]
  \item \textbf{Camera}: adaptive near/far ($0.005 \cdot d \dots 500 \cdot d$)
        keeps \texttt{depth24plus} precision usable across the whole zoom
        range; wheel-zoom band expanded to
        $10^{-3} \dots 10^4$ scene units
  \item \textbf{Camera target switching}: orbit-camera target switchable
        between central body / spacecraft / Moon via a footer dropdown
  \item \textbf{Apse markers}: tiny magenta/yellow crosses at periapsis /
        apoapsis. Detection via $\dot r$ zero-crossings (not triplet
        extremum) and suppressed when
        $(r_\text{max} - r_\text{min}) / \bar r < 10^{-3}$ to avoid spurious
        markers on near-circular orbits
  \item \textbf{Altitude chart}: $|\vec r|(t)$ polyline in the footer with
        a vertical cursor that tracks the time scrubber
\end{itemize}

\subsection{3.x — Known limitations}

\begin{itemize}[leftmargin=*]
  \item Multi-burn NLP returns inertial $\Delta v$; the editor speaks RIC,
        so it displays results rather than round-tripping. A future
        \texttt{dv\_inertial} field on the Maneuver model would close this.
  \item Moon ephemeris is sampled only at the configured demo epoch
        (\texttt{2026-01-01T00:00:00}); a per-trajectory configurable
        $t_0$ is a small follow-up.
  \item Frame is always inertial J2000. Phase~4 introduces the
        synodic-frame option needed for periodic-orbit analysis.
\end{itemize}

\subsection{3 — Status \& tests}

All 43 backend tests pass; ruff clean; frontend type-check clean.

% ============================================================
%  PHASE 4 — CR3BP, Lagrange points, low-energy transfers
% ============================================================
\newpage
\section{Phase 4 — Rotating Frames, CR3BP \& Low-Energy Transfers}

\subsection{Goals}

Add the dynamics, frames, and visualization needed to design and visualize
\emph{low-energy} (manifold-guided) transfers and \emph{ballistic capture}
trajectories. The Newtonian path remains the fast default; CR3BP is opt-in
via a frame selector and a separate propagator family.

The deliverable is: ``the user can pick an Earth--Moon halo orbit by energy,
see the unstable manifold tube reaching Earth, and choose an initial state
on that tube that ballistically captures into a low lunar orbit.''

\subsection{Pillar A --- Frame infrastructure}

\subsubsection{4.1 Frame definitions}

A \texttt{Frame} dataclass in \texttt{backend/oamp/frames.py} (new file)
with time-dependent rotation + translation:

\begin{itemize}[leftmargin=*]
  \item \textbf{J2000}: inertial (current default)
  \item \textbf{EM\_SYNODIC}: rotating with the Earth--Moon line, origin at
        the system barycenter --- the Moon and Earth are stationary
  \item \textbf{SE\_SYNODIC}: rotating with the Sun--Earth line; required
        for Weak-Stability-Boundary transfers
  \item \textbf{BCI\_EARTH} / \textbf{BCI\_MOON}: body-centered J2000
  \item \textbf{BCR\_*}: body-centered rotating, surface-relative
\end{itemize}

Each conversion is a $3{\times}3$ rotation $R(t)$ plus an origin
translation $\Delta\vec r(t)$, plus the velocity correction
$\vec v' = R(\vec v - \boldsymbol\omega \times \vec r) - \dot{\Delta\vec r}$.

\subsubsection{4.2 Frame-aware propagator and visualization}

\begin{itemize}[leftmargin=*]
  \item \texttt{PropagateRequest.output\_frame} --- backend transforms the
        result before returning
  \item \texttt{POST /transform/states} --- standalone batch transform
        endpoint
  \item Frontend frame selector dropdown coupled to \texttt{output\_frame}
        and to the renderer's ``what is at origin''
\end{itemize}

\subsection{Pillar B --- CR3BP dynamics}

\subsubsection{4.3 CR3BP propagator}

Non-dimensional equations of motion in the synodic frame with mass ratio
$\mu = m_2/(m_1+m_2)$.  In file
\texttt{backend/oamp/dynamics/cr3bp.py}.  JAX-friendly RHS so we can JIT
and \texttt{vmap} for batch trajectory studies (this links to the GPU
acceleration work):

\begin{equation}
  \begin{aligned}
    \ddot x &= \phantom{-}2\dot y + x - (1-\mu)\frac{x+\mu}{r_1^3} - \mu\frac{x-1+\mu}{r_2^3} \\
    \ddot y &= -2\dot x + y - (1-\mu)\frac{y}{r_1^3} - \mu\frac{y}{r_2^3} \\
    \ddot z &= - (1-\mu)\frac{z}{r_1^3} - \mu\frac{z}{r_2^3}
  \end{aligned}
\end{equation}

\subsubsection{4.4 Lagrange points + Jacobi constant}

\begin{itemize}[leftmargin=*]
  \item Newton solve for collinear $L_1, L_2, L_3$ from the quintic;
        exact for triangular $L_4, L_5$
  \item \texttt{jacobi\_constant(state, $\mu$)} as a pseudo-energy invariant;
        conserved to $\sim 10^{-10}$ over a closed periodic orbit, useful
        as an integration-error gauge
  \item Endpoint \texttt{GET /cr3bp/lagrange?mu=...}
\end{itemize}

\subsubsection{4.5 Periodic orbit families (hardest single item)}

\begin{itemize}[leftmargin=*]
  \item Planar Lyapunov around $L_1, L_2$ --- linearised initial guess
        + differential correction (Newton on the period-crossing residual)
  \item 3D halo orbits --- Richardson 3rd-order analytical seed, then
        differential correction
  \item Vertical Lyapunov (out-of-plane)
  \item \textbf{Family continuation}: pseudo-arclength on energy/period to
        walk along a family
  \item Endpoint
        \texttt{POST /cr3bp/periodic-orbit \{type, L, energy\}}
\end{itemize}

\subsubsection{4.6 Invariant manifolds}

\begin{itemize}[leftmargin=*]
  \item Monodromy matrix $\Phi(T)$ around a periodic orbit
  \item Eigendecomposition $\to$ stable ($|\lambda|<1$) and unstable
        ($|\lambda|>1$) directions
  \item Integrate the stable manifold backward and the unstable manifold
        forward from a sampling of points along the orbit
        ($\sim$100 trajectories typical)
  \item Endpoint
        \texttt{POST /cr3bp/manifold \{orbit, direction, branch, duration\}}
\end{itemize}

\subsection{Pillar C --- Low-energy transfers}

\subsubsection{4.7 Manifold-guided transfers}

\begin{itemize}[leftmargin=*]
  \item \textbf{Earth--Moon $L_1/L_2$ transit}: depart LEO $\to$ unstable
        manifold of $L_1/L_2$ halo $\to$ arrive at the lunar SOI without
        a capture burn
  \item \textbf{WSB (Belbruno) transfer}: depart Earth $\to$ Sun--Earth
        manifold tube $\to$ swing through Sun--Earth $L_1/L_2$ $\to$ arrive
        at Earth--Moon $L_2$ manifold $\to$ ballistic capture into lunar
        orbit. $\sim$90--110~d TOF, $\sim$10--20\% $\Delta v$ savings vs.\
        Hohmann (Hiten, GRAIL, LADEE precedent).
  \item Endpoint
        \texttt{POST /transfer/low-energy \{depart, target, type\}}
\end{itemize}

\subsubsection{4.8 Ballistic-capture diagnostic}

Sample many initial conditions near the lunar $L_1/L_2$ manifolds, propagate
backward in the BCR4BP (Sun perturbation included), classify each as
``captured'' (lunar bound for $\ge 3$ months) or ``escapes''. The boundary
between these classes is the \textbf{Weak Stability Boundary}.  Visualize as
a colour-coded map in the synodic frame --- the iconic Belbruno plot.

\subsection{Pillar C visualization}

\subsubsection{4.9 Frontend integration}

\begin{itemize}[leftmargin=*]
  \item Frame selector coupled to \texttt{output\_frame} on every propagate
        call; renderer ``central body'' becomes the barycenter when synodic
  \item $L_1$\,\dots\,$L_5$ markers (small $\times$ glyphs) drawn at their
        synodic-frame positions
  \item Periodic-orbit catalogue browser --- preset halo/Lyapunov entries
        by energy, click to load
  \item Manifold tube rendering: many semi-transparent trajectories,
        colour-coded by stable/unstable and $\pm$ branch
  \item Live Jacobi-constant readout in the footer (sanity check; drift
        diagnoses integration error or 4-body perturbation)
\end{itemize}

\subsection{Sequencing}

\begin{center}
\begin{tabular}{ll}
  \toprule
  4.1 \,\&\, 4.3 & foundational, independent, $\sim$2--3~d each \\
  4.2 (frame propagator) & needs 4.1; $\sim$1~d \\
  4.4 (Lagrange + Jacobi) & needs 4.3; $\sim$1~d \\
  4.5 (periodic orbits) & needs 4.4; $\sim$5--7~d (DC robustness) \\
  4.6 (manifolds) & needs 4.5; $\sim$3--4~d \\
  4.7 (transfers) & needs 4.6 + 4.2; $\sim$3~d \\
  4.8 (WSB diagnostic) & needs 4.6 + BCR4BP RHS; $\sim$4--5~d \\
  4.9 (visualization) & threads through all of the above; $\sim$3--4~d \\
  \midrule
  \textbf{Total} & \textbf{$\sim$5--6 weeks} \\
  \bottomrule
\end{tabular}
\end{center}

\subsection{Hard tradeoffs decided up front}

\begin{itemize}[leftmargin=*]
  \item \textbf{CR3BP vs full ephemeris}: do all family computation in
        pure CR3BP; refine into the SPICE-ephemeris model as a forward-
        shooting follow-up. Pure CR3BP first.
  \item \textbf{JAX vs SciPy for batch manifolds}: build on SciPy DOP853
        first to match the existing code path; expose a JAX
        \texttt{lax.scan} + \texttt{vmap} fast path under the \texttt{gpu}
        pixi feature in a follow-up.
  \item \textbf{Synodic-frame UI complexity}: do the inertial$\to$synodic
        transform on the backend (\texttt{output\_frame} in the request);
        renderer stays a dumb pipeline.
  \item \textbf{Discoverability}: add a small glossary section in the
        sidebar covering Jacobi, halo, manifold, WSB --- these terms are
        unusable to anyone outside the field.
  \item \textbf{BCR4BP scope}: only used by 4.8 (ballistic capture
        diagnostic). Don't try to do periodic orbits or generic
        propagation in BCR4BP --- CR3BP is enough everywhere else.
\end{itemize}

% ============================================================
%  PHASE 5  (originally drafted as Phase 3; renumbered after the actual
%  Phase 3 (sandbox integration + Moon) and Phase 4 (CR3BP) were inserted
%  ahead of it.  Content unchanged.)
% ============================================================
\newpage
\section{Phase 5 — Post-Newtonian \& General Relativity}

\subsection{Goals}

Provide GR-fidelity propagation for precision science missions: geodetic
surveying, pulsar timing, deep-space navigation, and any scenario where
Newtonian errors exceed the mission budget. The Newtonian path remains the
fast default; GR is opt-in via \texttt{physics: "pn1" | "gr"}.

\subsection{3a — Post-Newtonian Corrections (Bridge Layer)}

PN corrections are additive accelerations on top of the existing DOP853
integrator — no new integrator required. This makes validation straightforward:
the PN limit must reproduce known analytic results before the full GR engine
is trusted.

\subsubsection{Schwarzschild (1PN)}

The dominant solar-system GR effect. In isotropic coordinates:

\begin{equation}
  \mathbf{a}_\text{Schw} = \frac{\mu}{c^2 r^3}
  \left[
    \left(4\frac{\mu}{r} - v^2\right)\mathbf{r}
    + 4(\mathbf{r}\cdot\mathbf{v})\mathbf{v}
  \right]
\end{equation}

\subsubsection{Lense--Thirring (1PN Frame-Dragging)}

Spin--orbit coupling from Earth's angular momentum
$\mathbf{S} = I_\oplus \boldsymbol{\omega}_\oplus$:

\begin{equation}
  \mathbf{a}_\text{LT} = \frac{2G}{c^2 r^3}
  \left[
    \left(\frac{\mathbf{r}}{r}\times\mathbf{v}\right)\times\mathbf{S}
    + \frac{3}{r^2}\left(\mathbf{S}\cdot\mathbf{r}\right)
      \left(\mathbf{r}\times\mathbf{v}\right)
  \right]
\end{equation}

\subsubsection{de~Sitter / Geodetic Precession (1PN)}

Solar barycentric contribution to orbit precession:

\begin{equation}
  \boldsymbol{\Omega}_\text{dS} =
    \frac{3}{2}\frac{GM_\odot}{c^2 a_\oplus(1-e_\oplus^2)}
    \frac{\boldsymbol{n}_\oplus \times \mathbf{v}_\oplus}{|\mathbf{v}_\oplus|}
\end{equation}

\subsubsection{Validation Tests}

\begin{itemize}[noitemsep]
  \item Mercury perihelion precession: $42.98''$/century $\pm\,0.5\%$
  \item Gravity Probe~B geodetic precession: $6606.1\,$mas/yr $\pm\,1\%$
  \item Shapiro delay signature reproduced for Earth--Moon baseline
  \item PN accelerations vanish in the $c\to\infty$ limit (numerical check)
\end{itemize}

\subsection{3b — Full GR Geodesic Integrator}

\subsubsection{Metric}

DSX metric / Brumberg--Kopejkin Barycentric Celestial Reference System (BCRS).
For a test particle of negligible mass the line element is:

\begin{equation}
  ds^2 = -\left(1 - \frac{2w}{c^2}\right)c^2\,dt^2
         + \left(1 + \frac{2w}{c^2}\right)\delta_{ij}\,dx^i\,dx^j
         + \mathcal{O}(c^{-4})
\end{equation}

where $w = \sum_A G M_A / r_A$ is the Newtonian potential summed over all
Solar-system bodies $A$.

\subsubsection{State Vector}

\begin{equation}
  \mathbf{X} = \bigl(t_\text{TDB},\; x^\mu,\; u^\mu\bigr)
  \quad\in\quad \mathbb{R}^{1+4+4}
\end{equation}

with the constraint $g_{\mu\nu} u^\mu u^\nu = -c^2$ monitored as an invariant.

\subsubsection{Equations of Motion}

Geodesic equation with optional non-geodesic thrust $f^\alpha$:

\begin{equation}
  \frac{Du^\alpha}{d\tau}
  = \frac{du^\alpha}{d\tau}
    + \Gamma^\alpha{}_{\mu\nu}\, u^\mu u^\nu
  = f^\alpha
\end{equation}

Christoffel symbols $\Gamma^\alpha{}_{\mu\nu}$ are computed analytically from
the PN metric expansion (no numerical differentiation).

\subsubsection{Integrator}

DOP853 on the 8D geodesic state with adaptive step-size; fallback to a
symplectic Runge--Kutta--Nyström for long arcs. The norm constraint
$g_{\mu\nu}u^\mu u^\nu + c^2$ is logged as a quality indicator.

\subsubsection{Endpoint}

\texttt{POST /propagate/gr} — accepts BCRS initial conditions in SI plus
coordinate time span. Returns proper time $\tau$, coordinate trajectory
$x^\mu(\tau)$, and the 4-velocity $u^\mu(\tau)$.

\subsection{3c — Tetrad Thrust Model}

Converts coordinate-frame thrust commands to the local orthonormal tetrad:

\begin{enumerate}[noitemsep]
  \item Build tetrad $e^\mu{}_{(a)}$ from $u^\mu$ and the spacecraft attitude
  \item Decompose thrust in the tetrad frame:
        $f^\alpha = (f^{\hat{0}}, f^{\hat{1}}, f^{\hat{2}}, f^{\hat{3}})$
  \item Inject as right-hand side of the non-geodesic equation above
\end{enumerate}

\texttt{POST /propagate/gr} accepts \texttt{thrust\_tetrad\_N: [float, float, float]}
(body-frame thrust in Newtons; frame defined relative to $u^\mu$ and a
reference attitude quaternion).

\subsection{3d — Multi-body Metric Pre-bake (Ray Workers)}

\subsubsection{Motivation}

Evaluating the BCRS metric and its derivatives at every integration step for
every Solar-system body at full SPICE fidelity is the dominant cost of GR
propagation (${\sim}100\times$ slower than Newtonian). Pre-baking the metric
on a 4D grid eliminates repeated SPICE calls.

\subsubsection{Implementation}

\begin{enumerate}[noitemsep]
  \item A Ray actor pool (\texttt{metric\_build}) samples $w(\mathbf{r}, t)$
        and its spatial gradients on a Cartesian grid over the mission window
  \item Results stored as HDF5 in \texttt{data/metric\_cache/} with a
        SHA256 manifest (same pattern as SPICE kernels)
  \item The GR integrator interpolates bilinearly in space and linearly in time
        — typically ${\lesssim}100\,\mu$s per step vs ${\sim}10$~ms for live SPICE
\end{enumerate}

\begin{verbatim}
pixi run -e full build-metric \
  --start 2025-01-01 --end 2025-06-01 \
  --bodies earth,moon,sun,jupiter \
  --resolution-km 500
\end{verbatim}

\subsubsection{Cache Invalidation}

The manifest pins the SPICE kernel SHA256 hashes used during pre-bake. If any
kernel is updated the cache is considered stale and must be rebuilt.

\subsection{3e — Frontend}

\begin{center}
\begin{tabularx}{\textwidth}{L{3.5cm} X}
  \toprule
  \textbf{Feature} & \textbf{Detail} \\
  \midrule
  GR vs PN vs Newton &
    Three overlaid trajectories, colour-coded, rendered simultaneously.
    Toggle any layer on/off. \\[4pt]
  Residual panel &
    Side panel showing $|\mathbf{r}_\text{GR} - \mathbf{r}_\text{Newton}|$
    vs time; highlights when GR corrections exceed a user-set threshold. \\[4pt]
  Proper time display &
    Elapsed proper time $\tau$ alongside coordinate time $t$ in the HUD;
    shows the integrated gravitational + kinematic time dilation. \\[4pt]
  Metric visualiser &
    2D equatorial slice of the Newtonian potential $w(\mathbf{r})$ as a
    heatmap underlay in the 3D view; toggled via a button. \\
  \bottomrule
\end{tabularx}
\end{center}

\subsection{Phase 3 Sequencing}

\begin{center}
\begin{tabular}{L{5cm} C{2.5cm} L{5.5cm}}
  \toprule
  \textbf{Work item} & \textbf{Est.\ duration} & \textbf{Dependency} \\
  \midrule
  3a PN corrections        & 2 weeks    & Phase 2 \\
  3b GR geodesic integrator& 3--4 weeks & 3a \\
  3c Tetrad thrust model   & 1 week     & 3b \\
  3d Ray metric cache      & 2--3 weeks & 3b \\
  3e Frontend              & 2 weeks    & 3b (parallel with 3d) \\
  \midrule
  \textbf{Total}           & \textbf{10--12 weeks} & \\
  \bottomrule
\end{tabular}
\end{center}

% ============================================================
%  OVERALL TIMELINE
% ============================================================
\newpage
\section{Overall Timeline}

\subsection{Phase comparison}

\begin{center}
\rowcolors{2}{rowalt}{white}
\begin{tabularx}{\textwidth}{L{2.6cm} X X X X X}
  \rowcolor{tabhead}
    & \color{white}\textbf{P1}
    & \color{white}\textbf{P2}
    & \color{white}\textbf{P3}
    & \color{white}\textbf{P4}
    & \color{white}\textbf{P5} \\
  Physics
    & Newton + $J_2$
    & Newton + $J_2$\,\dots\,$J_6$ + drag + SRP + 3B
    & = P2 (no new physics)
    & CR3BP + BCR4BP
    & PN1 + full GR (BCRS) \\
  Frames
    & J2000
    & J2000
    & + body selector (Earth/Moon)
    & + EM/SE synodic + BCR
    & BCRS \\
  Integrator
    & DOP853
    & + Verlet + Yoshida-4
    & = P2
    & + symplectic SABA on CR3BP
    & + RKN for geodesics \\
  Optimisers
    & ---
    & Lambert + Hohmann + CasADi NLP
    & all four exposed in sidebar; + TLE ingest
    & + manifold-shooting + diff.\ correction
    & + GR-aware shooting \\
  Frontend
    & 3 demo scenes
    & Maneuver editor + scrubber
    & Solver panel, Moon render, cislunar demo, apse markers, altitude chart, camera targeting
    & Frame selector, $L_n$ markers, periodic-orbit catalogue, manifold tubes, WSB map
    & GR/PN/Newton compare, residual panel, metric heatmap \\
  Heavy compute
    & ---
    & Ray streaming
    & ---
    & JAX \texttt{vmap} for batch manifolds (\texttt{gpu} feature)
    & Ray (metric pre-bake) \\
  Status
    & \emph{shipped}
    & \emph{shipped}
    & \emph{shipped} (3.1--3.6)
    & planned, $\sim$5--6 weeks
    & planned, $\sim$10--12 weeks \\
\end{tabularx}
\end{center}

\subsection{Critical path}

\begin{center}
\begin{minipage}{0.85\textwidth}
\begin{verbatim}
   P1 ──► P2 ──► P3 ──┬──► P4 ──► low-energy + ballistic capture
                      │
                      └──► P5 ──► GR / precision science
\end{verbatim}
\end{minipage}
\end{center}

Phases 4 and 5 are independent of each other.  Phase 4 specialises the
dynamics towards three-body geometry (mission design, manifold transfers);
Phase 5 specialises towards relativistic effects (precision navigation,
geodesy).  They can be developed in parallel by different contributors once
the Phase 3 foundations are locked in.

% ============================================================
%  CONVENTIONS REFERENCE
% ============================================================
\section{Conventions Reference}

\begin{center}
\begin{tabularx}{\textwidth}{L{3cm} X}
  \toprule
  \textbf{Convention} & \textbf{Value} \\
  \midrule
  Licence        & Apache-2.0 (explicit patent grant) \\
  Time on wire   & UTC ISO-8601 string \\
  Time internal  & TDB seconds since J2000 (= SPICE ET); see \texttt{oamp.timescales} \\
  Units          & SI throughout the engine (m, m/s, kg, s) \\
  Accelerator    & JAX (pixi \texttt{gpu} feature); CPU default \\
  Heavy compute  & Ray (pixi \texttt{full} feature) \\
  SPICE kernels  & Pinned by SHA256 in \texttt{backend/oamp/spice/manifest.toml} \\
  Epoch J2000    & 2000-01-01T11:58:55.816Z UTC \\
  $TT - TAI$     & $32.184\,$s (fixed by IAU) \\
  \bottomrule
\end{tabularx}
\end{center}

\end{document}
